\documentclass{IOS-Book-Article}

\usepackage[T1]{fontenc}
\usepackage{mathptmx}
\usepackage{soul}\setuldepth{article}
\usepackage{longtable}
\usepackage{array}
\def\hb{\hbox to 11.5 cm{}}

\begin{document}

\pagestyle{headings}
\def\thepage{}
\begin{frontmatter}

\title{LexSR: Consolidating Legislation and Linking Case Law to Build\\
a Legal Knowledge Graph for Suriname, a Low-Resource Jurisdiction}

\markboth{}{2026\hb}

\author[A]{\fnms{Ernst} \snm{van Gassen}%
\thanks{Corresponding Author: Ernst van Gassen, E-mail: evg.gassen@gmail.com.}}

\runningauthor{van Gassen}
\address[A]{Independent researcher at Arktos AI Lab}

\begin{abstract}
Access to consolidated, machine-readable legislation and linked case law is a prerequisite
for effective legal information systems. No such resource previously existed for Suriname, a
Dutch-speaking jurisdiction of fewer than 600,000 people with a predominantly scanned legal
archive. We present LexSR, a pipeline that ingests raw government PDFs and
the Hof van Justitie's case-law portal, applies OCR, classifies documents, automatically
consolidates \emph{wijzigingswetten} (amendment acts) into versioned legislative expressions
modelled on FRBR, and links case law to the legislation it applies. The pipeline produces a
corpus of 1,623 works (1,323 legislative, 300 judicial) and 45,293 fragments, from
which we construct a unified knowledge graph of 2,768 edges spanning amendment,
citation, and case-application relationships. We report two evaluations: a synthetic
citation-resolution benchmark, where our resolver achieves 100\% precision/recall against a
curated 61-case gold set, and, as the paper's methodological contribution, a measurement of
how much of that resolution accuracy actually survives contact with the real, messily-titled
corpus it was built to run against. Only 17\% of extracted citations resolve to an
in-corpus work. We attribute this to a structural mismatch between the idealized canonical
identifiers assumed by curated benchmarks and the filename-derived identifiers that scraped,
OCR'd ingestion actually produces, and show that reconciling the two recovers most of the
shortfall. We also measure amendment-parsing coverage as the corpus grew
from 33 to 50 amendment documents (82\% $\rightarrow$ 38\%), a second instance of the same
divergence between curated-subset accuracy and real-corpus performance. We discuss
implications for low-resource legal NLP and make the corpus, code, and evaluation scripts
publicly available.
\end{abstract}

\begin{keyword}
legal knowledge graph\sep legislative consolidation\sep case-law citation linking\sep
low-resource legal NLP\sep FRBR\sep Suriname
\end{keyword}
\end{frontmatter}

\section{Introduction}

Consolidation (merging an amendment act's instructions into the base law it modifies) and
citation linkage (connecting case law to the provisions it applies) are both legally
critical and technically demanding. Most European jurisdictions maintain authoritative
consolidated legislation and case-law databases through governmental portals (e.g.,
\emph{wetten.overheid.nl} and \emph{rechtspraak.nl} in the Netherlands, EUR-Lex for EU law). Small and
developing jurisdictions often lack the institutional capacity for either, leaving
practitioners to manually track amendment chains and cross-reference case law with no
integrated view of the law in force.

Suriname illustrates this gap sharply. The \emph{Staatsblad} (official gazette) is available as
individual scanned PDFs through several government sources, and the Hof van Justitie
publishes case decisions on a dedicated portal (\emph{rechtspraak.sr}, launched 2019), but no
system links the two: a decision citing ``artikel 4 Arbeidswet 1963'' offers no way to jump to
that provision, and no portal shows any act's consolidated text as amended.

This paper makes four contributions:

\begin{enumerate}
\item \textbf{LexSR corpus}: a structured, FRBR-modelled database of 1,323 Surinamese legislative
   works and 300 Hof van Justitie decisions, with 45,293 fragments, derived from ten live
   government sources, the judicial case-law portal, and locally digitised archives.
\item \textbf{Amendment parser}: a hybrid regex + LLM system that extracts structured consolidation
   instructions from \emph{wijzigingswetten}, evaluated across the corpus's full growth trajectory.
\item \textbf{Case-law integration}: a scraper and ingestion pipeline for \emph{rechtspraak.sr} that
   segments decisions into numbered considerations and links them to cited legislation,
   reusing the same citation-resolution machinery built for legislation-to-legislation
   citations.
\item \textbf{Knowledge graph and an accuracy audit}: a unified graph of amendment, citation,
   and case-application edges over the full corpus, together with a measurement of the gap
   between synthetic-benchmark citation accuracy and real-corpus linkage rate, and an
   account of why the two diverge.
\end{enumerate}

The remainder of this paper is organised as follows. Section~\ref{sec:related} reviews related work.
Section~\ref{sec:corpus} describes the corpus and ingestion pipeline, including the newly added case-law
source. Section~\ref{sec:consolidation} details amendment consolidation. Section~\ref{sec:caselaw} describes case-law integration
and knowledge graph construction. Section~\ref{sec:eval} presents evaluation results, including the
citation-linkage accuracy audit. Section~\ref{sec:discussion} discusses limitations and future work, and
Section~\ref{sec:conclusion} concludes.

\section{Related Work}
\label{sec:related}

\subsection{Legal Document Modelling}

The FRBR (Functional Requirements for Bibliographic Records) model provides a
well-established abstraction for legislative documents, distinguishing the abstract Work
(a law), its Expression (a version at a given date), its Manifestation (a particular
encoding), and individual Items \cite{ifla1998}. Akoma Ntoso \cite{palmirani2011a} is
the primary XML schema for legislation structured around these concepts and has been adopted
by the UN, EU, and several national parliaments. The European Legislation Identifier (ELI)
\cite{pubsoffice2012} provides URI conventions compatible with this model. We extend
this modelling to case law following the same FRBR pattern: a court decision is a Work with
a single Expression (the published text does not get amended the way legislation does), and
its numbered considerations are Fragments, the case-law analogue of legislative articles.

\subsection{Legal Knowledge Systems for Consolidation}

Automated consolidation has received sustained attention in the legal knowledge systems
community. Biagioli, Francesconi et al. \cite{francesconi2004} describe the \emph{Norme in
Rete} project's XML standards and tools for cross-referencing and consolidating Italian
legislation. Spinosa et al. \cite{spinosa2009} apply NLP-based metadata extraction to detect
and classify textual amendments (repeal, substitution, integration) for the same
consolidation task. The \emph{legislation.gov.uk} system \cite{sheridan2010}
handles consolidation for UK law at scale with a semi-automated editorial workflow. Vitali
and Zeni \cite{vitali2007} discuss amendment formalisation in Akoma Ntoso. The Dutch
\emph{wetten.overheid.nl} portal maintains consolidated texts for all Dutch and former-colonial
legislation through a government editorial process, but does not publish its consolidation
tooling \cite{koop2022}. LexSR differs from these in targeting a low-resource setting (no
existing structured archive, predominantly scanned PDFs, and no editorial support) with a
fully automated pipeline.

\subsection{Legal Knowledge Graphs and Citation Networks}

Citation network analysis of case law has a long history in AI and Law, from early work on
precedent networks to more recent graph-based approaches to legal information retrieval and
argumentation mining \cite{benchcapon2012}. Legal knowledge graphs unifying legislation
and case law have been built for well-resourced jurisdictions; Dhani et al.
\cite{dhani2021}, for example, construct one from Indian court judgments and statutes for
case-similarity recommendation. These approaches rely on structured source data. We are
not aware of prior work constructing such a graph from scratch for a jurisdiction whose
primary source material is unstructured scanned PDFs and an unindexed WordPress-based case
portal, which is the setting this paper addresses.

\subsection{NLP for Legal Amendment Text}

Information extraction from legal amendment provisions has been studied for several European
languages. Lesmo, Mazzei, Palmirani et al. \cite{tulsi2013} present TULSI, an NLP system
combining deep syntactic analysis with a shallow semantic interpreter to extract and classify
modificatory provisions in Italian legislative texts. For Dutch, Boer et al. \cite{boer2007}
propose the MetaLex/LKIF XML standards for structuring legislation, and Hoekstra et al.
\cite{hoekstra2007} extend this with the LKIF Core ontology of basic legal concepts, but
explicit amendment/consolidation parsing for Dutch has received limited attention in the
JURIX and ICAIL literature. The use of large language models for
legal text understanding is by now standard \cite{chalkidis2020,koreeda2021}. Our hybrid
approach, regex-first with an LLM fallback for low-confidence cases, is motivated by the need for determinism and auditability while retaining flexibility for novel
constructions.

\subsection{Legal Corpora for Small Jurisdictions}

The JURIX and ICAIL communities have repeatedly called for legal NLP
resources beyond the major European languages \cite{niklaus2024}. Krasadakis et al.
\cite{krasadakis2024} survey the specific challenges legal informatics poses for
low-resource-language NLP. To our knowledge, no
machine-readable structured corpus linking legislation and case law for Suriname has
previously been published.

\section{Corpus and Ingestion Pipeline}
\label{sec:corpus}

\subsection{Sources}

LexSR ingests from ten live government sources plus locally digitised PDF archives: the
parliament website (DNA.sr), the legal database SRIS, the tax authority (Belastingdienst),
the central bank (MAS), the standards bureau (SSB), the main government portal, and four
ministry sites (finance, labor, natural resources, and the electoral bureau). We evaluated
all sixteen government sources our scraper infrastructure supports; six were excluded after
diagnosis showed their configured URLs no longer match the live site structure
(Section~\ref{sec:discussion} discusses this as a recurring low-resource-jurisdiction problem: source
sites change without notice and without a stable API to detect it). Case law is sourced separately from
\emph{rechtspraak.sr}, the Hof van Justitie's decision portal, described in Section~\ref{sec:caselaw-scraping}. All
source documents are stored immutably in a bronze layer with SHA-256 checksums and retrieval
timestamps.

\subsection{Data Model}

The canonical store implements FRBR at four levels, shown in Table~\ref{tab:frbr}.

\begin{table}[h]
\caption{FRBR levels and LexSR entities}
\label{tab:frbr}
\begin{tabular}{|l|l|l|l|}
\hline
\textbf{FRBR level} & \textbf{LexSR entity} & \textbf{Legislation example} & \textbf{Case-law example} \\
\hline
Work & \texttt{Work} & Arbeidswet 1963 & SRU-HvJ-2022-39 \\
Expression & \texttt{Expression} & version\_date=2024-01-15 & single expression \\
Manifestation & \texttt{Manifestation} & source PDF & source HTML page \\
Item & \texttt{Fragment} & Artikel 4, lid 2 & consideration 2.1 \\
\hline
\end{tabular}
\end{table}

Stable URIs follow an ELI-inspired pattern rooted at \texttt{/id/\{scheme\}/\{id\}}, with
\texttt{law} and \texttt{caselaw} as the two schemes currently populated. The amendment relation is
represented by an \texttt{AmendmentLink} table connecting an amending Work to the base Work it
modifies; the citation relation (legislation-to-legislation, case-to-legislation, and
case-to-case) is represented by a single generic \texttt{Citation} table typed by
\texttt{citation\_type}, avoiding a schema split between the two document kinds.

\subsection{Corpus Statistics}

Table~\ref{tab:stats} summarizes the corpus. The legislative corpus spans material from 1931 to 2025, with the densest coverage in the
2010 to 2025 period corresponding to the availability of digital-native PDFs. The case-law
sample spans 1952 to 2025 and was drawn from a discoverable pool of 900+ decisions on
\emph{rechtspraak.sr} (Section~\ref{sec:caselaw-scraping}); this paper reports on the first 300 ingested.

\begin{table}[h]
\caption{Corpus statistics}
\label{tab:stats}
\begin{tabular}{|l|r|}
\hline
\textbf{Metric} & \textbf{Value} \\
\hline
Legislative works & 1,323 \\
Case-law works (decisions) & 300 \\
Expressions & 1,633 \\
Fragments (articles, paragraphs, considerations) & 45,293 \\
Considerations extracted from case law & 1,142 \\
Citations extracted & 15,867 \\
Citations resolved to an in-corpus work & 2,713 (17.1\%) \\
Amendment links identified & 55 \\
Amendment links fully applied by the engine & 11 \\
Knowledge graph nodes / edges & 1,623 / 2,768 \\
\hline
\end{tabular}
\end{table}

\subsection{OCR and Structure Extraction}

For scanned PDFs, text is extracted using a pipeline that applies standard PDF text
extraction first and falls back to Tesseract OCR for image-based pages. Article boundaries
are detected heuristically using the standard Dutch/Surinamese legislative structure:
``Artikel N'' headings, ``Lid N'' sub-paragraphs, and lettered or numbered sub-items. OCR quality
varies substantially across the corpus, and a subset carries the lower text quality of an
English rather than a Dutch recognition model. We record this as a property of the released
corpus, since it bears on the parsing coverage reported in Section~\ref{sec:parsing-eval}.

\section{Amendment Consolidation}
\label{sec:consolidation}

\subsection{Structure of a \emph{Wijzigingswet}}

A Surinamese amendment act (\emph{wijzigingswet} or \emph{wijzigingsbesluit}) follows a standard
template. The operative section (Artikel 1 of the amendment) names the base law and lists
labelled instructions:

\begin{quote}
\small\ttfamily
Artikel 1\\
De Wet Toezicht en Controle op\\
Financi\"{e}le Dienstverlening (S.B. 2009 no. 78)\\
wordt als volgt gewijzigd:\\[4pt]
A. In artikel 3 lid 2 wordt "vijftien" vervangen door "tien".\\
B. Na artikel 16 wordt een nieuw artikel 16a\\
\ \ ingevoegd, luidende:\\
\ \ "Artikel 16a. ..."\\
C. Artikel 24 vervalt.
\end{quote}

The amendment act's later articles (Artikel 2, 3, ...) may amend a second base law or
specify the in-force date. The act typically concludes with a \emph{Memorie van Toelichting}
(explanatory memorandum), which must be excluded from instruction extraction.

\subsection{Parser Architecture}

The \texttt{AmendmentParser} operates in two passes. \textbf{Pass 1 (regex)} splits the document into
sections by ``Artikel N'' boundaries (after stripping explanatory-memorandum text), extracts
the target law name from the section header or preamble, splits labelled instruction blocks
(A., B., C. or Roman-numeral labels), and classifies each block against a priority-ordered
set of $\sim$20 regex patterns covering the standard Dutch amendment constructions (repeal,
replacement, insertion, phrase substitution). \textbf{Pass 2 (LLM fallback)}, used only when a
regex-classified instruction's confidence falls below 0.70, re-sends the section to an LLM
with a structured JSON prompt; results are merged, with LLM output overriding regex output
where its confidence is higher. The fallback targets any OpenAI-compatible endpoint, local
or hosted. It is disabled for every result reported in this paper, so all figures below
describe regex-only behaviour, consistent with how we measured at the smaller corpus size.

\subsection{Consolidation Engine}

\texttt{ConsolidationEngine} consumes \texttt{AmendmentInstruction} objects and applies them to the current
expression of the base law to produce a new \texttt{Expression} at the amendment's effective date.
Base-law resolution, which maps the amendment's raw target-law reference to a Work in the
store, uses fuzzy title matching (exact match, normalized match, then length-guarded
substring containment) rather than exact-ID lookup, because the amendment's stated law name
essentially never matches the filename-derived work ID that real ingestion produces. This
matching logic is shared with citation resolution, which requires the identical fallback for
the identical reason (Section~\ref{sec:linkage-gap}).

\section{Case-Law Integration and Knowledge Graph}
\label{sec:caselaw}

\subsection{Scraping \emph{rechtspraak.sr}}
\label{sec:caselaw-scraping}

\emph{rechtspraak.sr} is a WordPress-based site with no listing or search API; the only reliable
enumeration mechanism is the WordPress XML sitemap, which we parse recursively (sitemap
index $\rightarrow$ per-type sub-sitemaps $\rightarrow$ individual decision URLs). This alone surfaced 900+
decisions spanning 1952 to 2025. Each decision page exposes clean structured metadata (court,
case number, decision and publication dates, legal area) as label/value pairs, and full
decision text as HTML, considerably cleaner input than the OCR'd legislative PDFs, since
no scanning or recognition step is involved.

\subsection{Consideration Segmentation}

Surinamese court decisions structure their legal reasoning as numbered paragraphs under a
``De beoordeling'' (assessment) heading, using section.paragraph numbering (e.g., ``2.1'', ``4.7'')
rather than an explicit ``r.o.'' (\emph{rechtsoverweging}) label. We segment on this numbering
pattern directly; decisions without recognizable numbering (typically older or very short
judgments) are kept as a single fragment rather than discarded, so no decision's text is
lost to a structural-parsing failure. Across the 300 ingested decisions this produced 1,142
consideration-level fragments.

\subsection{Citation Extraction from Case Law}

Legislative citations within case-law text (``artikel 4 lid 2 Arbeidswet 1963'') are extracted
using the identical citation resolver and extraction pipeline built for legislation-to-
legislation citations. No case-law-specific resolution logic was needed, since a citation is
a citation regardless of source document type. This yielded 526 case-to-legislation
citations across the 300-decision sample. Case-to-case citations (a decision citing another
by its site-assigned identifier, e.g. ``SRU-HvJ-2019-3'') use a separate lightweight regex
match, since these identifiers are specific to this one source and not part of the general
citation-resolution grammar; none were found in this sample, which we discuss as an open
question in Section~\ref{sec:limitations} rather than assume is a parsing gap.

\subsection{Knowledge Graph Construction}

We assemble a unified graph directly from the canonical store, adding no schema or
storage layer, with one node per Work (typed \texttt{legislation} or \texttt{case\_law}) and three edge
types: AMENDS (from \texttt{AmendmentLink}), CITES (legislation-to-legislation citations), and
APPLIES (case-to-legislation citations). The resulting graph over the full corpus has
1,623 nodes and 2,768 edges (2,322 CITES, 391 APPLIES, 55 AMENDS), with a largest connected
component of 706 nodes. The ten most-cited works are, in order: Personeelswet (civil service
law), Burgerlijk Wetboek (Civil Code), Wetboek van Strafrecht (Criminal Code), Grondwet
(Constitution), Wetboek van Koophandel (Commercial Code), and several tax and electoral
statutes. This distribution matches expectation for a general legal system and offers a
face-validity check on the extraction pipeline independent of the benchmark in Section~\ref{sec:resolver-eval}.

\section{Evaluation}
\label{sec:eval}

\subsection{Citation Resolution: Synthetic Benchmark}
\label{sec:resolver-eval}

We evaluate the citation resolver on a curated 61-case known-answer dataset spanning
citation, temporal, and gazette-reference categories. The resolver achieves a 100\% pass
rate, 100\% precision/recall/F1, and 91.2\% raw resolution rate (the remainder being
deliberately-unresolvable negative test cases). Pattern coverage is therefore sound against
well-formed input, which is precisely the condition the next section removes.

\subsection{Citation Resolution: Real-Corpus Linkage Gap}
\label{sec:linkage-gap}

Synthetic-benchmark accuracy measures whether the resolver correctly identifies \emph{that} a
reference exists and what canonical URI it should map to. It does not measure whether that
URI corresponds to anything in the corpus. These are separable properties, and for LexSR
they diverged completely: every one of 13,949 citations extracted from the corpus carried a
null target work, despite each having passed resolution. The divergence is invisible to the
aggregate ``citations resolved'' statistic, which counts successful pattern matches rather
than successful linkage, and observable only in the \texttt{target\_work\_id} column. We
attribute it to three causes, of which the second is the one we consider generalizable.

\begingroup\sloppy
\begin{enumerate}
\item \textbf{Identifier-scheme drift.} 788 works and their expressions retained a superseded
   \texttt{sr/law/...} URI form from an earlier naming convention rather than \texttt{/id/law/...},
   so direct lookup could not match them.
\item \textbf{Idealized vs. filename-derived identifiers.} The resolver's abbreviation registry maps
   ``Grondwet'' to the canonical slug \texttt{grondwet-1987}, the identifier the corpus
   \emph{should} carry. Real ingestion, run against scraped and OCR'd documents with no
   naming authority, instead produces identifiers derived from source filenames
   (\texttt{grondwet-met-mvt-\allowbreak{}4fb9cced}). No amount of resolver accuracy against a synthetic
   benchmark addresses this, because the discrepancy is not in resolution logic; it is a
   mismatch between two independently reasonable naming schemes, one assumed by the
   benchmark and one produced by ingestion. Reconciling them requires resolution to fall
   back on fuzzy title matching, the same mechanism amendment base-law resolution already
   needed (Section~\ref{sec:consolidation}) for the same underlying reason.
\item \textbf{Provenance typing.} The citation-type classifier detects case-law provenance from
   an \texttt{/ecli/} URI substring, a convention this corpus does not use: its case-law
   scheme is \texttt{/id/caselaw/...} (Section~\ref{sec:corpus}). Case-law citations were
   therefore typed \texttt{unknown}.
\end{enumerate}
\endgroup

With all three reconciled, target-work resolution across the corpus rose from 0/13,949 to
2,713/15,867 (17.1\%; the citation count itself grew between measurements as case-law
ingestion added citations from decision text). The remaining 83\% are overwhelmingly gazette
references (\emph{Staatsblad}/\emph{Gouvernementsblad} citations), which are correctly extracted and
typed but cannot resolve further because the corpus does not model individual gazette issues
as first-class entities. This is a genuine scope boundary rather than a resolution failure,
and one we flag as future work (Section~\ref{sec:future}).

We report this gap because we believe it generalizes: a resolver validated only against a
synthetic or curated benchmark can report perfect accuracy while linking nothing in the
corpus it is meant to serve, if the benchmark's identifiers do not reflect how real
ingestion actually names things. We are not aware of this specific failure mode being
discussed in the low-resource legal NLP literature and suggest it as a checklist item for
similar projects: \emph{measure resolution against the real corpus's actual identifiers, not just
against curated ground truth.}

\subsection{Amendment Parsing: Coverage Over Corpus Growth}
\label{sec:parsing-eval}

We evaluate the amendment parser (regex pass only, matching the methodology we used when
only 33 amendment documents had been ingested) on all amendment documents currently in the corpus (a document may
target more than one base law, so this is evaluated per amendment document, not per
\texttt{AmendmentLink} row; the corpus has 55 links across 50 unique documents). Table~\ref{tab:parsing}
reports both measurements.

\begin{table}[h]
\caption{Amendment parsing coverage before and after corpus expansion}
\label{tab:parsing}
\begin{tabular}{|l|r|r|r|}
\hline
\textbf{Corpus state} & \textbf{Amendment documents} & \textbf{OK ($\geq$1 instruction)} & \textbf{Rate} \\
\hline
Original (curated subset) & 33 & 27 & 81.8\% \\
Current (full live-scraped corpus) & 50 & 19 & 38.0\% \\
\hline
\end{tabular}
\end{table}

The drop is a second instance of the same pattern as Section~\ref{sec:linkage-gap}: the original
33-document set was disproportionately drawn from cleanly-digitised, well-OCR'd sources, while
the expanded corpus includes material from ten additional live-scraped sources with more
variable OCR quality and, in some cases, genuinely different amendment-drafting
conventions we have not yet characterized. As at the smaller size, failures are
concentrated in cases where the operative instruction text itself is missing or garbled in
the source fragment, not in cases where legible operative text was misparsed. We have
not yet repeated the per-case failure classification (OCR garbling vs. missing
text vs. novel construction) at the new scale, and flag this as immediate future work
(Section~\ref{sec:future}) rather than assert a cause we have not verified at this corpus size.

\subsection{Engine Validation}

Eleven amendments have been fully applied by the consolidation engine (up from 2 at the
smaller corpus) and are available for inspection in the canonical store; the remainder of
the 19 parseable-but-unapplied cases fail engine-level base-law resolution or produce
low-confidence instructions withheld pending review, rather than being blocked by parsing.

\section{Discussion and Future Work}
\label{sec:discussion}

\subsection{Limitations}
\label{sec:limitations}

\textbf{Corpus coverage and source stability.} Six of sixteen candidate scraper sources were
excluded because their configured URLs no longer resolve or the site structure has changed.
This is a specific instance of a general low-resource-jurisdiction problem: government sites are
maintained inconsistently and rarely expose a stable API, so scraper configurations decay
silently over time with no error signal beyond ``zero documents found.'' Automated periodic
health-checking of scraper coverage is future infrastructure work, not yet built.

\textbf{Real-corpus linkage.} Section~\ref{sec:linkage-gap}'s 17.1\% real-corpus citation linkage rate, while a large
improvement on the unreconciled 0\%, remains low in absolute terms; the majority of the gap is
the un-modelled gazette-issue scope boundary rather than a resolution failure, but this
should be stated as a limitation of the current corpus's entity coverage, not treated as
solved.

\textbf{Amendment-parsing coverage.} The 38\% current-corpus OK rate (Section~\ref{sec:parsing-eval}) is a lower bar
than we would want for a production consolidation system, though it should be read alongside
the fact that ``OK'' here means at least one clean instruction was extracted per amendment,
and many practically important amendments are simpler, single-instruction acts well within
this capability.

\textbf{Case-to-case citation coverage.} Zero case-to-case citations were found in the 300-decision
sample (Section~\ref{sec:caselaw}); we do not yet know whether this reflects a genuine rarity of explicit
cross-citation in this corpus (plausible for a small jurisdiction with a correspondingly
small body of precedent) or an extraction-pattern gap, and consider this an open
methodological question rather than a settled negative result.

\subsection{Broader Implications}

Suriname shares its legislative drafting conventions (Dutch-language templates derived from
former colonial law) with the Dutch Caribbean islands (Aruba, Cura\c{c}ao, Sint Maarten). The
LexSR pipeline may be directly applicable to these jurisdictions with minimal adaptation.
More broadly, the paper's central methodological point, that synthetic benchmark accuracy
and real-corpus linkage rate are different measurements that can diverge sharply, applies
to any legal knowledge system built by combining a hand-curated resolution component with
uncurated, real-world ingestion, independent of jurisdiction or language.

\subsection{Future Work}
\label{sec:future}

\begin{enumerate}
\item \textbf{Failure-mode analysis at scale}: Repeat the per-case classification
   (OCR garbling vs. missing operative text vs. novel construction) across the current
   50-document amendment corpus.
\item \textbf{Gazette-issue modelling}: Represent individual \emph{Staatsblad}/\emph{Gouvernementsblad} issues
   as first-class entities to close the largest remaining share of the citation-linkage gap.
\item \textbf{Case-law scale-up}: Ingest the remaining 600+ decisions available on \emph{rechtspraak.sr}
   and re-evaluate case-to-case citation coverage at that scale.
\item \textbf{Gold-standard instruction annotation}: Annotate amendment acts at instruction level
   (target article, operation type, find/replace text) for precision/recall evaluation
   beyond the current binary OK/NO\_MATCH measure.
\item \textbf{Source health monitoring}: Automated detection of scraper configuration decay
   (Section~\ref{sec:limitations}) rather than manual discovery.
\item \textbf{User evaluation}: Study practitioner use of the portal and knowledge graph through a
   user study with Surinamese legal practitioners.
\end{enumerate}

\section{Conclusion}
\label{sec:conclusion}

We presented LexSR, a pipeline for building a structured, versioned corpus of Surinamese
legislation and case law, automatically consolidating amendment acts, and constructing a
unified knowledge graph linking legislation, amendments, and case-law citations. Beyond the
system itself, this paper's central contribution is methodological: we measured, rather than
assumed, how much of our citation resolver's synthetic-benchmark accuracy (100\% precision/
recall) survives contact with a real, uncurated corpus (17.1\% linkage), attributed the gap
to a structural mismatch between idealized and filename-derived identifiers, and showed that
reconciling the two recovers most of it.
A parallel measurement of amendment-parsing coverage across corpus growth (82\% $\rightarrow$ 38\%) tells
the same story from a different angle. We believe both findings generalize beyond this
jurisdiction: legal NLP systems evaluated only against curated benchmarks can silently fail
to serve the real corpora they are built for, and that gap is worth measuring directly. The
corpus (1,623 works, 45,293 fragments, 2,768 knowledge-graph edges), code, and evaluation
scripts are released under open licences to support future research on legal NLP for small
jurisdictions.

\begin{thebibliography}{99}

\bibitem{ifla1998}
IFLA Study Group. Functional Requirements for Bibliographic Records. IFLA Publications; 1998.

\bibitem{palmirani2011a}
Palmirani M, Vitali F. Akoma-Ntoso for Legal Documents. In: Sartor G, et al., editors.
Legislative XML for the Semantic Web. Springer; 2011. p. 75-100.

\bibitem{pubsoffice2012}
Publications Office of the European Union. European Legislation Identifier (ELI). Official
Journal 2012/C 325/02; 2012.

\bibitem{francesconi2004}
Biagioli C, Francesconi E, Spinosa P, Taddei M. XML Documents Within a Legal Domain:
Standards and Tools for the Italian Legislative Environment. In: Document Analysis Systems
VI (DAS 2004). Lecture Notes in Computer Science, vol 3163. Springer; 2004.

\bibitem{spinosa2009}
Spinosa P, Giardiello G, Cherubini M, Marchi S, Venturi G, Montemagni S. NLP-based metadata
extraction for legal text consolidation. In: Proceedings of the 12th International Conference
on Artificial Intelligence and Law (ICAIL 2009); 2009. p. 40-49.

\bibitem{sheridan2010}
Sheridan J, Tennison J. Linking UK government data. In: Proceedings of LDOW 2010; 2010.

\bibitem{vitali2007}
Vitali F, Zeni F. Towards a country-independent data format: the Akoma Ntoso experience. In:
Proceedings of the V Legislative XML Workshop. Florence; 2007.

\bibitem{koop2022}
KOOP. Offici\"{e}le bekendmakingen: wetten.overheid.nl. Kennis- en exploitatiecentrum voor
offici\"{e}le overheidspublicaties; 2022.

\bibitem{benchcapon2012}
Bench-Capon T, et al. A history of AI and Law in 50 papers. Artificial Intelligence and Law.
2012;20(3):215-319.

\bibitem{dhani2021}
Dhani JS, Bhatt R, Ganesan B, Sirohi P, Bhatnagar V. Similar Cases Recommendation using Legal
Knowledge Graphs. arXiv:2107.04771; 2021.

\bibitem{tulsi2013}
Lesmo L, Mazzei A, Palmirani M, Radicioni DP. TULSI: an NLP system for extracting legal
modificatory provisions. Artificial Intelligence and Law. 2013;21:139-172.

\bibitem{boer2007}
Boer A, Winkels R, Vitali F. Proposed XML standards for law: MetaLex and LKIF. In:
Proceedings of JURIX 2007. IOS Press; 2007.

\bibitem{hoekstra2007}
Hoekstra R, Breuker J, Di Bello M, Boer A. The LKIF core ontology of basic legal concepts. In:
Proceedings of the LOAIT Workshop. CEUR Workshop Proceedings, vol 321; 2007. p. 43-63.

\bibitem{chalkidis2020}
Chalkidis I, et al. LEGAL-BERT: The muppets straight out of law school. In: Findings of
EMNLP 2020; 2020.

\bibitem{koreeda2021}
Koreeda Y, Manning C. ContractNLI: A dataset for document-level natural language inference
for contracts. In: Findings of EMNLP 2021; 2021.

\bibitem{niklaus2024}
Niklaus J, et al. MultiLegalPile: A 689GB Multilingual Legal Corpus. In: Proceedings of ACL
2024 (Volume 1: Long Papers); 2024.

\bibitem{krasadakis2024}
Krasadakis P, Sakkopoulos E, Verykios VS. A Survey on Challenges and Advances in Natural
Language Processing with a Focus on Legal Informatics and Low-Resource Languages.
Electronics. 2024;13(3):648.

\end{thebibliography}
\end{document}
